% SEGMENT OLP-0113-S01
% Part: first-order-logic 
% Chapter: axiomatic-deduction 
% Section: rules-and-proofs

\documentclass[../../../include/open-logic-section]{subfiles}

\begin{document}

\iftag{FOL}
      {\olfileid{fol}{axd}{rul}}
      {\olfileid{pl}{axd}{rul}}

\olsection{விதிகளும் \usetoken{P}{derivation}}

% SEGMENT OLP-0113-S02
\begin{explain}
அடிகோள் சார்ந்த !!{derivation}s தருக்கவியலுக்கான மிக எளிய
!!{derivation} முறைமையாக இருக்கலாம். !!^a{derivation} என்பது
!!{formula}s இன் ஒரு தொடர் மட்டுமே. !!a{derivation} ஆகக் கருதப்பட,
தொடரிலுள்ள ஒவ்வொரு !!{formula}வும் ஓர் அடிகோளின் நிகழ்வாக இருக்க
வேண்டும்; அல்லது அதற்கு முன் தொடரில் வரும் ஒன்று அல்லது அதற்கு மேற்பட்ட
!!{formula}s இலிருந்து ஓர் அனுமான விதியால் பின்வர வேண்டும்.
!!^a{derivation} அதன் கடைசி !!{formula}வை !!{derive} செய்கிறது.
\end{explain}

% SEGMENT OLP-0113-S03
\begin{defn}[!!^{derivability}]
$\Gamma$ என்பது $\Lang L$ இன் !!{formula}s கணம் எனில்,
$\Gamma$ இலிருந்து ஒரு \emph{!!{derivation}} என்பது $!A_1$,
\dots,~$!A_n$ என்ற !!{formula}s இன் முடிவுறு தொடராகும்; ஒவ்வொரு
$i \le n$ க்கும் பின்வருவனவற்றில் ஒன்று பொருந்தும்:
\begin{enumerate}
\item $!A_i \in \Gamma$; அல்லது
\item $!A_i$ ஓர் அடிகோள்; அல்லது
\item $j < i$ (மற்றும் $k < i$) ஆகிய ஏதோ $!A_j$ (மற்றும் $!A_k$)
  இலிருந்து ஓர் அனுமான விதியால் $!A_i$ பின்வருகிறது.
\end{enumerate}
\end{defn}

% SEGMENT OLP-0113-S04
எது சரியான !!{derivation} ஆகக் கருதப்படுகிறது என்பது நாம் எந்த அனுமான
விதிகளை அனுமதிக்கிறோம் என்பதையும் (மேலும், எவற்றை அடிகோள்களாகக்
கொள்கிறோம் என்பதையும்) சார்ந்தது. சில நிபந்தனைகளின்கீழ்,
!!a{derivation} இல் உள்ள படி~$A_i$ சரியான அனுமானப் படி என்று கூறும்
எனில்--அப்போது கூற்றே ஓர் அனுமான விதி.

% SEGMENT OLP-0113-S05
\begin{defn}[அனுமான விதி]
$\Gamma$ இலிருந்து !!a{derivation} இல் எது சரியான அனுமானப் படியாகக்
கருதப்படுகிறது என்பதற்குப் போதுமான நிபந்தனையை ஓர் \emph{அனுமான விதி}
தருகிறது.
\end{defn}

% SEGMENT OLP-0113-S06
எடுத்துக்காட்டாக, $!A \in \Gamma$ எனும் ஓருறுப்புத் தொடர் $!A$ எதுவும்
அற்பமாகவே !!a{derivation} ஆகக் கருதப்படுவதால், பின்வருவது மிக எளிய
அனுமான விதியாக இருக்கலாம்:
\begin{quote}
  $!A \in \Gamma$ எனில், $\Gamma$ இலிருந்து எந்த !!{derivation} இலும்
  $!A$ எப்போதும் ஒரு சரியான அனுமானப் படி.
\end{quote}

% SEGMENT OLP-0113-S07
அதேபோல், $!A$ அடிகோள்களில் ஒன்று எனில், $!A$ மட்டுமே
!!a{derivation}; எனவே இதுவும் ஓர் அனுமான விதி:
\begin{quote}
  $!A$ ஓர் அடிகோள் எனில், $!A$ ஒரு சரியான அனுமானப் படி.
\end{quote}

% SEGMENT OLP-0113-S08
கருதப்படும் படிக்கு முன் வரும் !!{formula}s ஐ அனுமான விதி பயன்படுத்தும்
போது விவரம் மேலும் சுவையாகிறது. பின்வரும் விதி \emph{மோடஸ்
போனென்ஸ்} எனப்படுகிறது:
\begin{quote}
  $!B \lif !A$ மற்றும் $!B$ ஆகியவை !!{derivation} இல் மேலே வந்தால்,
  $!A$ ஒரு சரியான அனுமானப் படி.
\end{quote}

% SEGMENT OLP-0113-S09
இதுவே ஒரே அனுமான விதி எனில், மேலுள்ள !!{derivation} வரையறை பின்வருமாறு
அமைகிறது: ஒவ்வொரு $i \le n$ க்கும் பின்வருவனவற்றில் ஒன்று பொருந்தினால்,
அப்போதும் அப்போது மட்டுமே, $!A_1$, \dots,~$!A_n$ என்பது
!!a{derivation}:
\begin{enumerate}
\item $!A_i \in \Gamma$; அல்லது
\item $!A_i$ ஓர் அடிகோள்; அல்லது
\item ஏதோ $j < i$ க்கு $!A_j$ என்பது $!B \lif !A_i$ ஆகவும், ஏதோ
  $k < i$ க்கு $!A_k$ என்பது~$!B$ ஆகவும் உள்ளது.
\end{enumerate}
கடைசி உட்கூறு, $!A_i$ என்பது~$!A_j$ ($!B \lif !A_i$) மற்றும் $!A_k$
($!B$) ஆகியவற்றிலிருந்து மோடஸ் போனென்ஸ் விதியால் பின்வருகிறது என்கிறது.
$1$ இலிருந்து~$n$ வரை செல்ல முடிந்து, ஒவ்வொரு முறையும்~$\Gamma$ இல்
உள்ள, ஓர் அடிகோளாக உள்ள, அல்லது ஓர் அனுமான விதி சரியான அனுமானப் படி எனக்
கூறும் !!a{formula}~$!A_i$ ஐக் கண்டால், முழுத் தொடரும் சரியான
!!{derivation} ஆகக் கருதப்படும்.

% SEGMENT OLP-0113-S10
\begin{defn}[!!^{derivability}]
$\Gamma$ இலிருந்து தொடங்கி $!A$ இல் முடியும் !!a{derivation} இருந்தால்,
!!a{formula}~$!A$ ஆனது $\Gamma$ இலிருந்து \emph{!!{derivable}};
இதனை $\Gamma \Proves !A$ என எழுதுகிறோம்.
\end{defn}

% SEGMENT OLP-0113-S11
\begin{defn}[தேற்றங்கள்]
வெற்றுக் கணத்திலிருந்து~$!A$ இன் !!a{derivation} இருந்தால்,
!!^a{formula}~$!A$ ஒரு \emph{தேற்றம்}. $!A$ ஒரு தேற்றம் எனில்
$\Proves !A$ என்றும், இல்லையெனில் $\Proves/ !A$ என்றும் எழுதுகிறோம்.
\end{defn}


\end{document}
